\chapter{Riemann Integration}
\label{ch:ii10}

Riemann integration quantifies accumulation by controlled finite sums. The theory rests on upper and lower sums, oscillation, and the principle that sufficiently fine partitions suppress local variation.

\section*{Learning goals}
After this chapter, the reader should be able to:
\begin{itemize}
\item compute upper and lower sums for explicit partitions and bounded functions;
\item prove Riemann integrability by controlling \(U(f,P)-L(f,P)\);
\item evaluate elementary Riemann integrals from partitions, antiderivatives, or comparison arguments;
\item show that monotone or continuous functions on compact intervals are Riemann integrable;
\item construct bounded functions that fail to be Riemann integrable and identify the obstruction;
\item use additivity and comparison properties of the integral to derive estimates;
\end{itemize}
\section*{Conceptual roadmap}
\[
\boxed{\text{definition}\;\longrightarrow\;\text{estimate}\;\longrightarrow\;\text{theorem}\;\longrightarrow\;\text{application}.}
\]

\section{Partitions and sums}
Tagged partitions produce Riemann sums; upper and lower sums eliminate dependence on tags.

\section{Riemann integrability}
A bounded function is integrable when upper and lower integrals coincide.

\section{Darboux criterion}
Integrability is equivalent to making the gap between upper and lower sums arbitrarily small.


% BEGIN VOL02-EXPANSION II10-example-01
\begin{example}[A Darboux-gap proof for a Lipschitz function]\label{ex:ii10-expand-01}
Suppose \(f\) is \(L\)-Lipschitz on \([a,b]\). On a partition interval of
length \(\Delta x_i\), its oscillation is at most \(L\Delta x_i\). Therefore
\[
U(f,P)-L(f,P)
\le L\sum_i(\Delta x_i)^2
\le L\,\|P\|\sum_i\Delta x_i
=L(b-a)\|P\|.
\]
Choosing mesh
\[
\|P\|<\frac{\varepsilon}{L(b-a)}
\]
makes the Darboux gap below \(\varepsilon\), proving integrability with an
explicit quantitative criterion.
\end{example}
% END VOL02-EXPANSION II10-example-01
\section{Continuous functions}
Uniform continuity on a compact interval makes continuous functions Riemann integrable.


% BEGIN VOL02-EXPANSION II10-example-02
\begin{example}[Uniform continuity drives continuous integrability]\label{ex:ii10-expand-02}
Let \(f\) be continuous on \([a,b]\). Given \(\varepsilon>0\), uniform
continuity supplies \(\delta>0\) such that intervals shorter than \(\delta\)
have oscillation below \(\varepsilon/(b-a)\). Choose a partition with mesh
below \(\delta\). Then
\[
U-L\le\frac{\varepsilon}{b-a}\sum_i\Delta x_i=\varepsilon.
\]
The proof reveals exactly where compactness enters: it upgrades pointwise
continuity to one uniform spatial scale.
\end{example}
% END VOL02-EXPANSION II10-example-02
\section{Monotone functions}
Monotonicity controls total oscillation and also guarantees integrability.

\section{Algebra of integrable functions}
Sums, scalar multiples, products, and absolute values preserve integrability under standard hypotheses.

\section{Fundamental theorem}
Integration and differentiation are inverse operations for continuous integrands.


% BEGIN VOL02-EXPANSION II10-example-03
\begin{example}[Fundamental theorem checked by a difference quotient]\label{ex:ii10-expand-03}
Let
\[
F(x)=\int_0^x t^2\,dt.
\]
Then
\[
\frac{F(x+h)-F(x)}h
=\frac1h\int_x^{x+h}t^2\,dt.
\]
For \(h>0\), the average lies between the minimum and maximum of \(t^2\) on the
shrinking interval; both tend to \(x^2\) as \(h\to0\). The same holds for
\(h<0\). Therefore \(F'(x)=x^2\), illustrating the local-average mechanism of
the fundamental theorem.
\end{example}
% END VOL02-EXPANSION II10-example-03
\section{Improper limits preview}
Classical Riemann integration on compact intervals is the base case for later extension to unbounded domains or singular integrands.

\section{Core structural results}

\begin{theorem}[Darboux criterion]\label{thm:ii10-01}
A bounded function on $[a,b]$ is Riemann integrable iff for every $\varepsilon>0$ there is a partition $P$ with $U(f,P)-L(f,P)<\varepsilon$.
\begin{proof}
One direction is immediate from equality of upper and lower integrals. Conversely a partition with small gap traps all finer upper and lower sums within epsilon, forcing the two integrals to coincide.
\end{proof}
\end{theorem}

\begin{theorem}[Continuous integrability]\label{thm:ii10-02}
Every continuous function on a compact interval is Riemann integrable.
\begin{proof}
Uniform continuity makes oscillation on sufficiently short subintervals uniformly small; summing oscillation times interval length controls the Darboux gap.
\end{proof}
\end{theorem}

\begin{theorem}[Fundamental theorem, Part I]\label{thm:ii10-03}
If $f$ is continuous on $[a,b]$ and $F(x)=\int_a^x f(t)dt$, then $F'(x)=f(x)$.
\begin{proof}
The difference quotient is an average value of $f$ over a shrinking interval; continuity makes that average converge to $f(x)$.
\end{proof}
\end{theorem}

\section{Worked examples}

\begin{example}[Integral of a step function]\label{ex:ii10-worked-01}
Compute the integral of $f=2$ on $[0,1)$ and $f=5$ on $[1,3]$. The integral is area sum $2\cdot1+5\cdot2=12$.
\end{example}

\begin{example}[Darboux gap estimate]\label{ex:ii10-worked-02}
If oscillation of $f$ on every partition interval is at most $\eta$, bound $U-L$. The gap is at most $\eta\sum\Delta x_i=\eta(b-a)$.
\end{example}

\begin{example}[Continuous integrability mechanism]\label{ex:ii10-worked-03}
Explain where compactness enters the proof. Continuity on $[a,b]$ becomes uniform continuity, so one partition scale works simultaneously across the whole interval.
\end{example}

\begin{example}[Monotone integrability]\label{ex:ii10-worked-04}
Show an increasing bounded function on $[a,b]$ is integrable. For an equal partition into $n$ pieces, the upper-lower gap telescopes to $(b-a)(f(b)-f(a))/n$.
\end{example}

\section{Solved dossiers}
The dossiers are longer solved problems. Legacy mappings guide topic selection where explicit retained source blocks exist; otherwise fresh canonical problems complete the chapter.

\begin{problem}[Integral of a step function]\label{prob:ii10-dossier-01}
Compute the integral of $f=2$ on $[0,1)$ and $f=5$ on $[1,3]$.
\end{problem}
\begin{solution}
The integral is area sum $2\cdot1+5\cdot2=12$.
\end{solution}

\begin{problem}[Darboux gap estimate]\label{prob:ii10-dossier-02}
If oscillation of $f$ on every partition interval is at most $\eta$, bound $U-L$.
\end{problem}
\begin{solution}
The gap is at most $\eta\sum\Delta x_i=\eta(b-a)$.
\end{solution}

\begin{problem}[Continuous integrability mechanism]\label{prob:ii10-dossier-03}
Explain where compactness enters the proof.
\end{problem}
\begin{solution}
Continuity on $[a,b]$ becomes uniform continuity, so one partition scale works simultaneously across the whole interval.
\end{solution}

\begin{problem}[Monotone integrability]\label{prob:ii10-dossier-04}
Show an increasing bounded function on $[a,b]$ is integrable.
\end{problem}
\begin{solution}
For an equal partition into $n$ pieces, the upper-lower gap telescopes to $(b-a)(f(b)-f(a))/n$.
\end{solution}

\begin{problem}[Integral of $x$]\label{prob:ii10-dossier-05}
Compute $\int_0^1x\,dx$ from equal Riemann sums.
\end{problem}
\begin{solution}
Right-endpoint sums are $n^{-2}\sum_{k=1}^nk=(n+1)/(2n)\to1/2$.
\end{solution}

\begin{problem}[A single discontinuity]\label{prob:ii10-dossier-06}
Why is a bounded function with one jump discontinuity still Riemann integrable?
\end{problem}
\begin{solution}
Isolate the jump in an interval of arbitrarily small length; on the remaining compact pieces continuity or small oscillation controls the Darboux gap.
\end{solution}

\begin{problem}[Dirichlet function]\label{prob:ii10-dossier-07}
Why is $1_{\mathbb Q}$ not Riemann integrable on any nontrivial interval?
\end{problem}
\begin{solution}
Every subinterval contains rationals and irrationals, so every upper sum is the interval length and every lower sum is zero.
\end{solution}

\begin{problem}[Absolute-value inequality]\label{prob:ii10-dossier-08}
Prove $|\int f|\le\int|f|$ for Riemann-integrable $f$.
\end{problem}
\begin{solution}
Pointwise $-|f|\le f\le|f|$; monotonicity of the integral gives the estimate.
\end{solution}

\begin{problem}[Mean value for integrals]\label{prob:ii10-dossier-09}
If $f$ is continuous on $[a,b]$, show $\int_a^bf=f(c)(b-a)$ for some $c$.
\end{problem}
\begin{solution}
The integral average lies between the minimum and maximum of $f$; the intermediate-value theorem gives a point attaining that average.
\end{solution}

\begin{problem}[FTC computation]\label{prob:ii10-dossier-10}
Evaluate $\int_0^1 3x^2\,dx$.
\end{problem}
\begin{solution}
An antiderivative is $x^3$, so the integral is $1$.
\end{solution}

\begin{problem}[Additivity]\label{prob:ii10-dossier-11}
Explain why $\int_a^cf+\int_c^bf=\int_a^bf$.
\end{problem}
\begin{solution}
Combine partitions on the two subintervals or compare Riemann sums; the common endpoint has no effect on the sum.
\end{solution}

\begin{problem}[Integration synthesis]\label{prob:ii10-dossier-12}
Why is Riemann integration a limit of finite-dimensional data?
\end{problem}
\begin{solution}
Each sum uses finitely many sample values and interval lengths; integrability asserts that all sufficiently fine such finite approximations converge to one common number.
\end{solution}

\section{Exercises with complete solutions}

\begin{exercise}\label{exr:ii10-01}
Compute $\int_0^2 1\,dx$.
\end{exercise}
\begin{hint}
Area of a rectangle. Write the upper and lower sums from the suprema and infima on each subinterval; keep the partition lengths explicit.
\end{hint}
\begin{solution}
It is $2$.
\end{solution}

\begin{exercise}\label{exr:ii10-02}
Compute $\int_0^1x^2\,dx$.
\end{exercise}
\begin{hint}
Use FTC. To prove integrability, aim directly at \(U(f,P)-L(f,P)<\varepsilon\) rather than at the value of the integral.
\end{hint}
\begin{solution}
It is $1/3$.
\end{solution}

\begin{exercise}\label{exr:ii10-03}
Is every continuous function on $[a,b]$ integrable?
\end{exercise}
\begin{hint}
Use compactness. For a monotone function, bound the oscillation on each subinterval by endpoint differences and sum the resulting telescoping estimate.
\end{hint}
\begin{solution}
Yes.
\end{solution}

\begin{exercise}\label{exr:ii10-04}
Is every bounded function integrable?
\end{exercise}
\begin{hint}
Recall Dirichlet's example. For a continuous function on a compact interval, use uniform continuity to make every subinterval oscillation small.
\end{hint}
\begin{solution}
No.
\end{solution}

\begin{exercise}\label{exr:ii10-05}
State additivity over adjacent intervals.
\end{exercise}
\begin{hint}
Split at $c$. Use additivity over adjacent intervals to reduce a complicated integral to simpler pieces.
\end{hint}
\begin{solution}
$\int_a^b=\int_a^c+\int_c^b$.
\end{solution}

\begin{exercise}\label{exr:ii10-06}
If $f\ge0$, what can be said about the integral?
\end{exercise}
\begin{hint}
Use monotonicity. Comparison follows from comparing upper or lower sums term by term.
\end{hint}
\begin{solution}
The integral is nonnegative.
\end{solution}

\begin{exercise}\label{exr:ii10-07}
Does changing one function value change the Riemann integral?
\end{exercise}
\begin{hint}
Think of partition sums. To show nonintegrability, prove that every partition leaves a fixed positive gap between upper and lower sums.
\end{hint}
\begin{solution}
No, for an integrable function a finite modification does not change the integral.
\end{solution}

\begin{exercise}\label{exr:ii10-08}
What is $\int_a^a f$?
\end{exercise}
\begin{hint}
Zero interval length. For a tagged Riemann-sum computation, separate the mesh-size estimate from the sample-point choice.
\end{hint}
\begin{solution}
It is $0$.
\end{solution}


% BEGIN VOL02-EXPANSION II10-exercises-01
\section{Graded supplementary exercises}
These exercises extend the preserved foundational set and are graded by mathematical role.

\subsection*{Standard computations}

\begin{exercise}[Step integral]\label{exr:ii10-expand-01}
Integrate the function equal to \(2\) on \([0,1)\) and \(5\) on \([1,2]\).
\end{exercise}
\begin{hint}
Sum rectangle areas.
\end{hint}
\begin{solution}
The integral is \(2\cdot1+5\cdot1=7\).
\end{solution}

\begin{exercise}[Upper lower sums]\label{exr:ii10-expand-02}
For \(f(x)=x\) on \([0,1]\) with partition \(\{0,1/2,1\}\), compute Darboux upper and lower sums.
\end{exercise}
\begin{hint}
Use monotonicity on each subinterval.
\end{hint}
\begin{solution}
\(L=0(1/2)+(1/2)(1/2)=1/4\); \(U=(1/2)(1/2)+1(1/2)=3/4\).
\end{solution}

\begin{exercise}[Integral bound]\label{exr:ii10-expand-03}
If \(|f|\le3\) on \([1,5]\), bound \(|\int_1^5f|\).
\end{exercise}
\begin{hint}
Use interval length times sup bound.
\end{hint}
\begin{solution}
At most \(4\cdot3=12\).
\end{solution}

\begin{exercise}[Monotone gap]\label{exr:ii10-expand-04}
For increasing \(f\) and equal \(n\)-partition of \([a,b]\), what Darboux-gap bound arises?
\end{exercise}
\begin{hint}
The differences telescope.
\end{hint}
\begin{solution}
\((b-a)(f(b)-f(a))/n\).
\end{solution}

\begin{exercise}[FTC]\label{exr:ii10-expand-05}
If \(F(x)=\int_0^x\cos t\,dt\), compute \(F'(x)\).
\end{exercise}
\begin{hint}
Use the fundamental theorem.
\end{hint}
\begin{solution}
\(F'(x)=\cos x\).
\end{solution}

\subsection*{Proofs}

\begin{exercise}[Darboux refinement]\label{exr:ii10-expand-06}
Explain why refining a partition cannot increase the lower sum or decrease the upper sum.
\end{exercise}
\begin{hint}
On smaller intervals, infima can only rise and suprema can only fall.
\end{hint}
\begin{solution}
Summing the refined contributions gives \(L(P)\le L(P')\le U(P')\le U(P)\).
\end{solution}

\begin{exercise}[Continuous integrability]\label{exr:ii10-expand-07}
Prove continuous functions on compact intervals are Riemann integrable.
\end{exercise}
\begin{hint}
Use uniform continuity to control oscillations on a fine partition.
\end{hint}
\begin{solution}
Choose mesh small enough that each oscillation is below \(\varepsilon/(b-a)\), making \(U-L<\varepsilon\).
\end{solution}

\begin{exercise}[Monotone integrability]\label{exr:ii10-expand-08}
Prove a monotone bounded function on \([a,b]\) is Riemann integrable.
\end{exercise}
\begin{hint}
Use equal partitions and telescope the oscillation sum.
\end{hint}
\begin{solution}
The Darboux gap is bounded by \((b-a)(f(b)-f(a))/n\to0\).
\end{solution}

\begin{exercise}[Integral comparison]\label{exr:ii10-expand-09}
If \(f\le g\) and both are integrable, prove \(\int f\le\int g\).
\end{exercise}
\begin{hint}
Compare lower/upper sums or integrate \(g-f\ge0\).
\end{hint}
\begin{solution}
Nonnegativity of \(g-f\) gives \(\int(g-f)\ge0\), hence the result.
\end{solution}

\subsection*{Counterexamples and hypothesis testing}

\begin{exercise}[Dirichlet function]\label{exr:ii10-expand-10}
Why is the indicator of \(\mathbb Q\cap[0,1]\) not Riemann integrable?
\end{exercise}
\begin{hint}
Every interval contains rational and irrational points.
\end{hint}
\begin{solution}
Every lower sum is \(0\) and every upper sum is \(1\), so the Darboux gap never shrinks.
\end{solution}

\begin{exercise}[Bounded insufficient]\label{exr:ii10-expand-11}
Does boundedness alone imply Riemann integrability?
\end{exercise}
\begin{hint}
Use the rational indicator.
\end{hint}
\begin{solution}
No; the Dirichlet function is bounded but nowhere locally low-oscillation.
\end{solution}

\begin{exercise}[Unbounded classical case]\label{exr:ii10-expand-12}
Why is \(1/\sqrt{x}\) on \((0,1]\) not covered by the classical bounded Riemann theory on \([0,1]\)?
\end{exercise}
\begin{hint}
Check boundedness near zero.
\end{hint}
\begin{solution}
It is unbounded at \(0\); its integral is treated as an improper limit, not a classical bounded Riemann integral on the closed interval.
\end{solution}

\subsection*{Applications and investigations}

\begin{exercise}[Numerical quadrature error intuition]\label{exr:ii10-expand-13}
Why does small oscillation on partition cells support accurate Riemann-sum quadrature?
\end{exercise}
\begin{hint}
Upper and lower sums squeeze all tagged sums.
\end{hint}
\begin{solution}
If \(U-L\) is small, every tagged sum lies in a narrow interval around the integral.
\end{solution}

\begin{exercise}[Accumulated error]\label{exr:ii10-expand-14}
If a uniform integrand approximation has error at most \(\eta\) on \([a,b]\), bound integral error.
\end{exercise}
\begin{hint}
Integrate the absolute difference.
\end{hint}
\begin{solution}
The error is at most \((b-a)\eta\).
\end{solution}

\subsection*{Challenge problems}

\begin{exercise}[Lipschitz integrability rate]\label{exr:ii10-expand-15}
Derive a mesh bound guaranteeing Darboux gap below \(\varepsilon\) for an \(L\)-Lipschitz function.
\end{exercise}
\begin{hint}
Oscillation on a cell is at most \(L\Delta x_i\).
\end{hint}
\begin{solution}
\(U-L\le L(b-a)\|P\|\), so take \(\|P\|<\varepsilon/[L(b-a)]\).
\end{solution}

\begin{exercise}[Zero integral nonnegative function]\label{exr:ii10-expand-16}
If continuous \(f\ge0\) on \([a,b]\) and \(\int f=0\), prove \(f\equiv0\).
\end{exercise}
\begin{hint}
If \(f(x_0)>0\), continuity gives a neighborhood with a positive lower bound.
\end{hint}
\begin{solution}
That neighborhood would contribute positive area, contradicting zero integral.
\end{solution}
% END VOL02-EXPANSION II10-exercises-01
\section*{Chapter summary}
The chapter's tools now form part of the canonical Volume II analysis toolkit and will be used in later approximation, fixed-point, and convergence arguments.
